\labsection{3}{Supervised learning II: model families}{sec:lab03-supervised-models}

\begin{labproject}{Neighbours, naive Bayes, trees, ensembles, and support vector machines on one workflow}
\textbf{Goal.} Compare model families under Lab~2's fixed validation workflow; inspect boundaries, tree depth, ensemble gains, and SVM tuning.

\medskip\noindent\textbf{Setup.}\par
\begin{lstlisting}[style=mlpython]
from pathlib import Path

import joblib
import matplotlib.pyplot as plt
import numpy as np
import pandas as pd
from sklearn.compose import ColumnTransformer
from sklearn.ensemble import GradientBoostingRegressor, RandomForestRegressor
from sklearn.impute import SimpleImputer
from sklearn.inspection import DecisionBoundaryDisplay, permutation_importance
from sklearn.linear_model import LogisticRegression
from sklearn.metrics import accuracy_score, mean_absolute_error, mean_squared_error, roc_auc_score
from sklearn.model_selection import (
    GridSearchCV, KFold, StratifiedKFold, cross_val_score, train_test_split,
)
from sklearn.naive_bayes import GaussianNB
from sklearn.neighbors import KNeighborsClassifier
from sklearn.pipeline import Pipeline
from sklearn.preprocessing import OneHotEncoder, StandardScaler
from sklearn.svm import SVC
from sklearn.tree import DecisionTreeClassifier, export_text, plot_tree

PROJECT_ROOT = Path.cwd().parent
DATA = PROJECT_ROOT / "all_datasets"
MODELS = PROJECT_ROOT / "models"
MODELS.mkdir(exist_ok=True)
\end{lstlisting}

\medskip\noindent\textbf{Part A --- KNN versus Gaussian naive Bayes.}\par

KNN votes among standardized neighbors; Gaussian naive Bayes combines per-feature Gaussian likelihoods and class priors.

\begin{lstlisting}[style=mlpython]
cancer = pd.read_csv(DATA / "breast_cancer_dataset.csv")

X_c = cancer.drop(columns=["id", "diagnosis", "Unnamed: 32"], errors="ignore")
y_c = (cancer["diagnosis"] == "M").astype(int)

Xc_train, Xc_valid, yc_train, yc_valid = train_test_split(
    X_c, y_c, test_size=0.25, random_state=42, stratify=y_c
)
\end{lstlisting}

$k$ is a hyperparameter, so it is chosen by cross-validation inside the training partition. Small $k$ follows every local quirk (low bias, high variance); large $k$ smooths everything (high bias, low variance).

\begin{lstlisting}[style=mlpython]
cv = StratifiedKFold(n_splits=5, shuffle=True, random_state=42)
k_values = [1, 3, 5, 7, 11, 21]
cv_scores = {}

for k in k_values:
    candidate = Pipeline([
        ("scale", StandardScaler()),
        ("model", KNeighborsClassifier(n_neighbors=k, weights="distance")),
    ])
    scores = cross_val_score(candidate, Xc_train, yc_train, cv=cv, scoring="roc_auc")
    cv_scores[k] = scores.mean()
    print(f"k={k:2d}  cross-validated ROC AUC = {scores.mean():.3f}")

best_k = max(cv_scores, key=cv_scores.get)
print("selected k:", best_k)

plt.figure(figsize=(5, 3.2))
plt.plot(k_values, [cv_scores[k] for k in k_values], marker="o")
plt.axvline(best_k, linestyle="--", color="grey")
plt.xlabel("k (number of neighbours)")
plt.ylabel("Cross-validated ROC AUC")
plt.title("Choosing k on training data only")
plt.tight_layout()
plt.show()
\end{lstlisting}

\textbf{Inspect.} Validation improves beyond $k=1$ and is similar for $k=7$--21.

With $k$ locked, both models are compared on the validation split.

\begin{lstlisting}[style=mlpython]
models = {
    "KNN": Pipeline([
        ("scale", StandardScaler()),
        ("model", KNeighborsClassifier(n_neighbors=best_k, weights="distance")),
    ]),
    "GaussianNB": Pipeline([
        ("scale", StandardScaler()),
        ("model", GaussianNB()),
    ]),
}

for name, model in models.items():
    model.fit(Xc_train, yc_train)
    pred = model.predict(Xc_valid)
    prob_c = model.predict_proba(Xc_valid)[:, 1]
    print(f"{name:10s} accuracy={accuracy_score(yc_valid, pred):.3f} "
          f"roc_auc={roc_auc_score(yc_valid, prob_c):.3f}")
\end{lstlisting}

Compare two-feature boundaries on tumor radius and texture, adding logistic regression as a reference.

\begin{lstlisting}[style=mlpython]
two = ["radius_mean", "texture_mean"]
fig, axes = plt.subplots(1, 3, figsize=(12, 3.8))
for ax, (name, estimator) in zip(axes, [
    ("Logistic regression", LogisticRegression(max_iter=2000)),
    (f"KNN (k={best_k})", KNeighborsClassifier(n_neighbors=best_k, weights="distance")),
    ("Gaussian naive Bayes", GaussianNB()),
]):
    small = Pipeline([("scale", StandardScaler()), ("model", estimator)])
    small.fit(Xc_train[two], yc_train)
    DecisionBoundaryDisplay.from_estimator(small, Xc_valid[two], ax=ax, alpha=0.3,
                                           response_method="predict", cmap="coolwarm")
    ax.scatter(Xc_valid[two[0]], Xc_valid[two[1]], c=yc_valid, cmap="coolwarm",
               s=12, edgecolor="k", linewidth=0.3)
    ax.set_title(f"{name}: {accuracy_score(yc_valid, small.predict(Xc_valid[two])):.2f} accuracy")
    ax.set_xlabel("radius_mean")
    ax.set_ylabel("texture_mean")
plt.tight_layout()
plt.show()
\end{lstlisting}

\textbf{Inspect.} Logistic regression is linear; Gaussian naive Bayes curves; KNN follows local votes.

\medskip\noindent\textbf{Part B --- One readable decision tree, and what happens as it grows.}\par

Trees split on value order, so numerical preprocessing imputes missingness without scaling.

\begin{lstlisting}[style=mlpython]
titanic = pd.read_csv(DATA / "titanic_dataset.csv")

features = ["Pclass", "Sex", "Age", "SibSp", "Parch", "Fare", "Embarked"]
t_num = ["Age", "SibSp", "Parch", "Fare"]
t_cat = ["Pclass", "Sex", "Embarked"]
X_t = titanic[features]
y_t = titanic["Survived"]

tree_preprocess = ColumnTransformer([
    ("num", SimpleImputer(strategy="median"), t_num),
    ("cat", Pipeline([
        ("impute", SimpleImputer(strategy="most_frequent")),
        ("onehot", OneHotEncoder(handle_unknown="ignore")),
    ]), t_cat),
])

Xt_train, Xt_valid, yt_train, yt_valid = train_test_split(
    X_t, y_t, test_size=0.25, random_state=42, stratify=y_t
)

tree = Pipeline([
    ("preprocess", tree_preprocess),
    ("tree", DecisionTreeClassifier(max_depth=4, min_samples_leaf=12, random_state=42)),
])
tree.fit(Xt_train, yt_train)
print("training accuracy:", round(accuracy_score(yt_train, tree.predict(Xt_train)), 3))
print("validation accuracy:", round(accuracy_score(yt_valid, tree.predict(Xt_valid)), 3))
\end{lstlisting}

Print and plot the tree; inspect split questions, sample counts, and classes.

\begin{lstlisting}[style=mlpython]
feature_names = list(tree.named_steps["preprocess"].get_feature_names_out())
print(export_text(tree.named_steps["tree"], feature_names=feature_names, max_depth=3))

plt.figure(figsize=(12, 6))
plot_tree(tree.named_steps["tree"], feature_names=feature_names,
          class_names=["did not survive", "survived"], max_depth=2, filled=True, fontsize=8)
plt.title("Top of the fitted Titanic decision tree")
plt.tight_layout()
plt.show()
\end{lstlisting}

\textbf{Inspect.} The tree first splits on sex, then age for men and class for women.

Fit each depth and compare training with validation accuracy.

\begin{lstlisting}[style=mlpython]
depths = list(range(1, 13))
train_acc, valid_acc = [], []
for depth in depths:
    t = Pipeline([
        ("preprocess", tree_preprocess),
        ("tree", DecisionTreeClassifier(max_depth=depth, random_state=42)),
    ]).fit(Xt_train, yt_train)
    train_acc.append(accuracy_score(yt_train, t.predict(Xt_train)))
    valid_acc.append(accuracy_score(yt_valid, t.predict(Xt_valid)))

plt.figure(figsize=(6, 3.6))
plt.plot(depths, train_acc, marker="o", label="training")
plt.plot(depths, valid_acc, marker="o", label="validation")
plt.xlabel("max_depth")
plt.ylabel("Accuracy")
plt.title("Deeper trees memorise; validation accuracy stops improving")
plt.legend()
plt.tight_layout()
plt.show()
\end{lstlisting}

\textbf{Inspect.} Training accuracy rises with depth while validation peaks earlier, exposing overfitting.

\medskip\noindent\textbf{Part C --- Random forest versus gradient boosting.}\par

Model $\log(1+\text{price})$ from Ames Housing's mixed-type features. Compare random forest and gradient boosting on identical three-fold splits.

\begin{lstlisting}[style=mlpython]
ames = pd.read_csv(DATA / "ames_housing_dataset.csv")
X_a = ames.drop(columns=["SalePrice", "Order", "PID"], errors="ignore")   # identifiers out
y_a = np.log1p(ames["SalePrice"])

ames_preprocess = ColumnTransformer([
    ("num", SimpleImputer(strategy="median"), X_a.select_dtypes(include="number").columns),
    ("cat", Pipeline([
        ("impute", SimpleImputer(strategy="most_frequent")),
        ("onehot", OneHotEncoder(handle_unknown="ignore", sparse_output=False)),
    ]), X_a.select_dtypes(exclude="number").columns),
])
ensembles = {
    "random_forest": RandomForestRegressor(
        n_estimators=150, max_features=0.8, min_samples_leaf=2, random_state=42, n_jobs=-1
    ),
    "gradient_boosting": GradientBoostingRegressor(
        n_estimators=150, learning_rate=0.05, max_depth=3, loss="huber", random_state=42
    ),
}

cv3 = KFold(n_splits=3, shuffle=True, random_state=42)
cv_rmse = {}
for name, estimator in ensembles.items():
    pipe = Pipeline([("preprocess", ames_preprocess), ("model", estimator)])
    neg_mse = cross_val_score(pipe, X_a, y_a, cv=cv3, scoring="neg_mean_squared_error")
    cv_rmse[name] = np.sqrt(-neg_mse)
    print(f"{name:18s} CV log-RMSE = {cv_rmse[name].mean():.4f} +/- {cv_rmse[name].std():.4f}")

plt.figure(figsize=(5, 3.4))
for i, (name, values) in enumerate(cv_rmse.items()):
    plt.scatter([i] * len(values), values, s=60, label="fold scores" if i == 0 else None)
    plt.hlines(values.mean(), i - 0.2, i + 0.2, color="black", label="mean" if i == 0 else None)
plt.xticks(range(len(cv_rmse)), list(cv_rmse))
plt.ylabel("log-RMSE per fold (lower is better)")
plt.title("Fold-by-fold: is the gap bigger than the spread?")
plt.legend()
plt.tight_layout()
plt.show()
\end{lstlisting}

\textbf{Inspect.} Boosting's observed gain is modest relative to fold variation. Weigh it against tuning and runtime.

Repeat on eleven wine measurements. Permute validation columns to measure the fitted forest’s reliance on each.

\begin{lstlisting}[style=mlpython]
wine = pd.read_csv(DATA / "red_wine_quality_dataset.csv")
X_w = wine.drop(columns="quality")
y_w = wine["quality"]
Xw_train, Xw_valid, yw_train, yw_valid = train_test_split(X_w, y_w, test_size=0.25, random_state=42)

wine_models = {
    "random_forest": RandomForestRegressor(n_estimators=400, min_samples_leaf=2,
                                           random_state=42, n_jobs=-1),
    "gradient_boosting": GradientBoostingRegressor(random_state=42),
}
wine_pred = {}
for name, model in wine_models.items():
    model.fit(Xw_train, yw_train)
    wine_pred[name] = model.predict(Xw_valid)
    print(f"{name:18s} MAE = {mean_absolute_error(yw_valid, wine_pred[name]):.3f}  "
          f"RMSE = {np.sqrt(mean_squared_error(yw_valid, wine_pred[name])):.3f}")

importance = permutation_importance(
    wine_models["random_forest"], Xw_valid, yw_valid, n_repeats=10, random_state=42, n_jobs=-1
)
ranked = pd.Series(importance.importances_mean, index=X_w.columns).sort_values()
spread = importance.importances_std[np.argsort(importance.importances_mean)]

fig, axes = plt.subplots(1, 2, figsize=(11, 3.8))
axes[0].barh(ranked.index, ranked.values, xerr=spread)
axes[0].set_xlabel("Increase in error when the column is shuffled")
axes[0].set_title("Permutation importance (random forest)")
axes[1].scatter(yw_valid, wine_pred["gradient_boosting"], alpha=0.5)
axes[1].plot([3, 8], [3, 8], linestyle="--", color="grey")
axes[1].set_xlabel("Observed quality")
axes[1].set_ylabel("Predicted quality")
axes[1].set_title("Gradient boosting pulls extremes to the middle")
plt.tight_layout()
plt.show()
\end{lstlisting}

\textbf{Inspect.} Alcohol and sulphates rank highly; predictions compress rare extreme ratings. Forests beat boosting here, unlike Ames.

\medskip\noindent\textbf{Part D --- Support vector machines: two dials and a kernel.}\par

Fit SVM scaling within each training fold. $C$ controls violation penalties; the kernel controls boundary shape; RBF $\gamma$ controls locality.

\begin{lstlisting}[style=mlpython]
X_train, X_test, y_train, y_test = train_test_split(
    X_c, y_c, test_size=0.25, random_state=42, stratify=y_c
)

svm_pipe = Pipeline([("scale", StandardScaler()), ("model", SVC())])
search = GridSearchCV(
    svm_pipe,
    param_grid={
        "model__kernel": ["linear", "rbf"],
        "model__C": [0.1, 1, 10, 100],
        "model__gamma": ["scale", 0.01, 0.1],
    },
    cv=StratifiedKFold(5, shuffle=True, random_state=42),
    scoring="roc_auc",
    n_jobs=-1,
)
search.fit(X_train, y_train)

score = search.decision_function(X_test)
pred = search.predict(X_test)
print("best parameters:", search.best_params_)
print("cross-validated ROC AUC:", round(search.best_score_, 3))
print("test accuracy:", round(accuracy_score(y_test, pred), 3))
print("test ROC AUC:", round(roc_auc_score(y_test, score), 3))
\end{lstlisting}

Plot cross-validated AUC for each $(C,\gamma)$ pair and compare two-feature kernel boundaries.

\begin{lstlisting}[style=mlpython]
results = pd.DataFrame(search.cv_results_)
rbf = results[results["param_model__kernel"] == "rbf"].copy()
rbf["gamma"] = rbf["param_model__gamma"].astype(str)
grid = rbf.pivot(index="param_model__C", columns="gamma", values="mean_test_score")

fig, axes = plt.subplots(1, 3, figsize=(13, 3.8))
im = axes[0].imshow(grid.values, cmap="viridis", aspect="auto")
axes[0].set_xticks(range(len(grid.columns)), grid.columns)
axes[0].set_yticks(range(len(grid.index)), grid.index)
axes[0].set_xlabel("gamma")
axes[0].set_ylabel("C")
for i in range(grid.shape[0]):
    for j in range(grid.shape[1]):
        axes[0].text(j, i, f"{grid.values[i, j]:.3f}", ha="center", va="center",
                     color="white", fontsize=8)
axes[0].set_title("RBF kernel: CV ROC AUC by C and gamma")
for ax, (title, estimator) in zip(axes[1:], [
    ("Linear kernel", SVC(kernel="linear", C=1)),
    ("RBF kernel (gamma = 1)", SVC(kernel="rbf", C=1, gamma=1.0)),
]):
    small = Pipeline([("scale", StandardScaler()), ("model", estimator)]).fit(X_train[two], y_train)
    DecisionBoundaryDisplay.from_estimator(small, X_test[two], ax=ax, alpha=0.3,
                                           response_method="predict", cmap="coolwarm")
    ax.scatter(X_test[two[0]], X_test[two[1]], c=y_test, cmap="coolwarm",
               s=12, edgecolor="k", linewidth=0.3)
    ax.set_title(f"{title}: {accuracy_score(y_test, small.predict(X_test[two])):.2f} accuracy")
    ax.set_xlabel("radius_mean")
    ax.set_ylabel("texture_mean")
plt.tight_layout()
plt.show()
\end{lstlisting}

\textbf{Inspect.} Most scores are similar. RBF allows curved boundaries, but these scaled classes gain little over a linear model.

Save, reload, and compare predictions with the in-memory model before sharing the artifact.

\begin{lstlisting}[style=mlpython]
model_path = MODELS / "breast_cancer_svm.joblib"
joblib.dump(search.best_estimator_, model_path)
loaded = joblib.load(model_path)
same = np.array_equal(loaded.predict(X_test), pred)
print("file size (KB):", round(model_path.stat().st_size / 1024, 1))
print("reloaded model reproduces the test predictions:", same)
\end{lstlisting}

\begin{tryit}
Grow the Titanic tree with \texttt{max\_depth=None} and \texttt{min\_samples\_leaf=1}, and compare training and validation accuracy with the depth-4 tree. Then refit the SVM search with only the linear kernel and compare its test ROC AUC with the RBF result.
\end{tryit}
\end{labproject}
